% Chapter 1

\chapter{Introducción general} % Main chapter title

\label{Chapter1} % For referencing the chapter elsewhere, use \ref{Chapter1} 
\label{IntroGeneral}

Todos los capítulos deben comenzar con un breve párrafo introductorio que indique cuál es el contenido que se encontrará al leerlo.  La redacción sobre el contenido de la memoria debe hacerse en presente y todo lo referido al proyecto en pasado, siempre de modo impersonal.

%----------------------------------------------------------------------------------------

% Define some commands to keep the formatting separated from the content 
\newcommand{\keyword}[1]{\textbf{#1}}
\newcommand{\tabhead}[1]{\textbf{#1}}
\newcommand{\code}[1]{\texttt{#1}}
\newcommand{\file}[1]{\texttt{\bfseries#1}}
\newcommand{\option}[1]{\texttt{\itshape#1}}
\newcommand{\grados}{$^{\circ}$}

%----------------------------------------------------------------------------------------
%\section{Introducción}

%----------------------------------------------------------------------------------------
\section{Estado del arte}

Revisión de tecnologías actuales para redes de sensores en invernaderos (cableadas e inalámbricas), comparando alternativas como Zigbee, LoRa, Wi-Fi y Bluetooth Mesh, con sus ventajas y limitaciones.



%----------------------------------------------------------------------------------------

\section{Motivación}

Explicación del problema (limitaciones del sistema cableado de Wentux Tecnoagro), la necesidad de escalabilidad y flexibilidad, y la relevancia de implementar una solución embebida basada en IoT.

%----------------------------------------------------------------------------------------

\section{Alcance y objetivos}

Definición de lo que cubre el trabajo y lo que queda fuera del alcance. Objetivo general y objetivos específicos.